\cotefile{157-4}
\scanpage{115}
\mspage{423}
% archive sheet number: 114
% source lines 1-13
C'était là sans doute une tâche urgente. J'ai fini par comprendre pourtant qu'il y avait plus urgent encore. Il m'a fallu d'abord parer au plus pressé : renouer le contact rompu avec mon corps, l'aider à remonter de l'état d'épuisement que j'ai fini par sentir et par admettre, et à retrouver la vigueur disparue. J'ai compris alors que pour cela, il fallait que je renonce pour une durée indéterminée à toute activité intellectuelle - fut-ce celle de méditer sur le sens de ce qui m'arrivait. C'est avec les notes reprises aujourd'hui que prend fin cette longue et salutaire ``parenthèse'' dans mes grands investissements, lesquels pour un temps (\unsure{depuis} le mois de février de cette année) s'étaient rejoints dans l'écriture de ``R\,é\,c\,o\,l\,t\,e\,s\quad e\,t\quad S\,e\,m\,a\,i\,l\,l\,e\,s''. La présente note est une toute première réflexion, ou tout au moins une sorte de compte-rendu sommaire, au sujet de cette ``parenthèse'' de quatre mois.

% source lines 14-21
Le temps de comprendre, à la fin des fins, le besoin d'un repos complet, une grande fatigue était devenue un épuisement profond. Faute d'avoir su écouter le langage pourtant péremptoire de mon corps, les quelques dérisoires pages de commentaires et retouches à l'Enterrement, arrachées sur un état de fatigue physique en ces premières deux semaines, l'ont été au prix d'une mise d'énergie qui, avec le recul, me paraît démentielle ! Toujours est-il qu'après ces prouesses, il m'a fallu rester allongé pendant de longues semaines, ne me levant que quelques heures par jour pour les tâches pratiques indispensables.

% source lines 22-28
Chose remarquable, une fois c\,o\,m\,p\,r\,i\,s le besoin d'un repos complet, je n'ai pas éprouvé la moindre difficulté à décrocher complètement de toute activité intellectuelle, sans aucune velléité de ``tricher''. Je n'avais pas même à prendre de décision à proprement parler - du seul fait d'avoir compris, j'avais déjà décroché. Les tâches qui la veille encore m'avaient tenu en haleine, soudain semblaient très lointaines, comme appartenant à un passé très reculé...

% source lines 29-36
Le présent n'en était pas vide pour autant. Alors que pendant des semaines et des mois le sommeil est resté réticent à venir, et que je restais allongé de longues heures, apparemment dans l'inaction totale, je ne me rappelle pas une seule fois avoir trouvé le temps long. Je refaisais connaissance avec mon corps, et aussi avec l'environnement le plus immédiat - ma chambre, ou parfois le morceau de gazon ou d'herbes sèches baigné de soleil juste devant mes yeux, là où d'aventure je m'étais allongé, près de la maison ou au cours d'une courte (et prudente...) promenade. Je passais de longs moments à poursuivre la danse d'une
